\section{Sammelproben}

\subsection{Der Befehl \texttt{\textbackslash sammleprobe}}

Der Befehl \verb|\sammelprobe| erzeugt einen standardisierten Kasten für Sammelproben.

Er enthält automatisch die Regeln:

\begin{itemize}
\item 3 TaW* entsprechen 1 QS
\item 10 QS sind notwendig
\item Teilerfolg bei 6 QS
\end{itemize}

Die angegebenen Talente und die Dauer werden automatisch formatiert.

\subsubsection*{Syntax}

\begin{Verbatim}
\sammelprobe
    {Talent(e)}
    {Dauer}
    [Ziel]
\end{Verbatim}

Der dritte Parameter ist optional.

\subsubsection*{Parameter}

\begin{center}
\begin{tabular}{|l|l|}
\hline
Parameter & Bedeutung\\
\hline
Talent(e) & verwendete Talente oder Fähigkeiten\\
\hline
Dauer & Zeit pro Versuch\\
\hline
Ziel & optionales Ziel der Probe\\
\hline
\end{tabular}
\end{center}

\subsubsection*{Beispiel mit Ziel}

\paragraph{Quellcode}

\begin{Verbatim}
\sammelprobe
    {Wildnisleben}
    {30 Minuten}
    [Einen sicheren Lagerplatz finden]
\end{Verbatim}

\paragraph{Ergebnis}

\sammelprobe
    {Wildnisleben}
    {30 Minuten}
    [Einen sicheren Lagerplatz finden]


\subsubsection*{Beispiel ohne Ziel}

\paragraph{Quellcode}

\begin{Verbatim}
\sammelprobe
    {Orientierung}
    {1 Stunde}
\end{Verbatim}

\paragraph{Ergebnis}

\sammelprobe
    {Orientierung}
    {1 Stunde}
